\section{Estimación de máxima verosimilitud y otros métodos}

\begin{ejemplo}
    Sea $X_1, \ldots, X_n$ una m.a.s. de $X \rightsquigarrow \cc{N}(\mu, \sigma^2)$, $\mu \in \mathbb{R}$, $\sigma^2 > 0$. Encontrar el EMV para $\mu$ y $\sigma^2$, en el caso de un parámetro conocido y cuando ambos parámetros son desconocidos. Además, deducir que el EMV no tiene por qué ser insesgado.
    \begin{enumerate}
        \item \ul{$\sigma_0^2$ conocido:} La función de verosimilitud es
        \begin{align*}
            L_{x_1, \ldots, x_n}(\mu) &= \prod_{i=1}^{n} f_{X_i}(x_i; \mu, \sigma_0^2) = \prod_{i=1}^{n} \frac{1}{\sqrt{2\pi\sigma_0^2}} \exp(-\frac{(x_i - \mu)^2}{2\sigma_0^2}) =\\&= \left(\frac{1}{\sqrt{2\pi\sigma_0^2}}\right)^n \exp\left(-\frac{1}{2\sigma_0^2} \sum_{i=1}^{n} (x_i - \mu)^2\right)\qquad \forall \mu \in \mathbb{R}.
        \end{align*}

        Puesto que el logaritmo es una función estrictamente creciente, podemos trabajar con la log-verosimilitud:
        \begin{align*}
            \ln(L_{x_1, \ldots, x_n}(\mu)) &= -\frac{n}{2}\ln(2\pi\sigma_0^2) - \frac{1}{2\sigma_0^2} \sum_{i=1}^{n} (x_i - \mu)^2\qquad \forall \mu \in \mathbb{R}.
        \end{align*}

        Como buscamos el máximo, derivamos e igualamos a cero:
        \begin{align*}
            \frac{d}{d\mu} \ln(L_{x_1, \ldots, x_n}(\mu)) &= -\frac{1}{2\sigma_0^2} \sum_{i=1}^{n} 2(x_i - \mu)(-1) = \frac{1}{\sigma_0^2} \sum_{i=1}^{n} (x_i - \mu) = 0 \implies \\&\implies \sum_{i=1}^{n} x_i - n\mu = 0 \implies \mu = \frac{1}{n} \sum_{i=1}^{n} x_i
        \end{align*}

        Sería necesario comprobar que efectivamente se trata de un máximo, pero como la función es cuadrática con coeficiente principal negativo, podemos concluir que efectivamente se trata de un máximo. Por lo tanto, el candidato a EMV para $\mu$ es
        \begin{align*}
            \hat{\mu} = \frac{1}{n} \sum_{i=1}^{n} X_i = \ol{X}.
        \end{align*}

        Sería necesario comprobar que este estadístico efectivamente es un estimador de $\mu$. Como $\mu\in \mathbb{R}$ y $\hat{\mu}\in \mathbb{R}$, podemos concluir que efectivamente $\hat{\mu}$ es un estimador de $\mu$, y por lo tanto, $\hat{\mu}$ es el EMV de $\mu$.

        \item \ul{$\mu_0$ conocido:} La función de verosimilitud es
        \begin{align*}
            L_{x_1, \ldots, x_n}(\sigma^2) &= \prod_{i=1}^{n} f_{X_i}(x_i; \mu_0, \sigma^2) = \prod_{i=1}^{n} \frac{1}{\sqrt{2\pi\sigma^2}} \exp(-\frac{(x_i - \mu_0)^2}{2\sigma^2}) =\\&= \left(\frac{1}{\sqrt{2\pi\sigma^2}}\right)^n \exp\left(-\frac{1}{2\sigma^2} \sum_{i=1}^{n} (x_i - \mu_0)^2\right)\qquad \forall \sigma^2 > 0.
        \end{align*}

        Puesto que el logaritmo es una función estrictamente creciente, podemos trabajar con la log-verosimilitud:
        \begin{align*}
            \ln(L_{x_1, \ldots, x_n}(\sigma^2)) &= -\frac{n}{2}\ln(2\pi\sigma^2) - \frac{1}{2\sigma^2} \sum_{i=1}^{n} (x_i - \mu_0)^2\qquad \forall \sigma^2 > 0.
        \end{align*}

        Como buscamos el máximo, derivamos e igualamos a cero:
        \begin{align*}
            \frac{d}{d\sigma^2} \ln(L_{x_1, \ldots, x_n}(\sigma^2)) &= -\frac{n}{2\sigma^2} + \frac{1}{2(\sigma^2)^2} \sum_{i=1}^{n} (x_i - \mu_0)^2 = 0 \implies \\&\implies -n\sigma^2 + \sum_{i=1}^{n} (x_i - \mu_0)^2 = 0 \implies \sigma^2 = \frac{1}{n} \sum_{i=1}^{n} (x_i - \mu_0)^2
        \end{align*}

        Se puede comprobar que efectivamente se trata de un máximo. Por lo tanto, el candidato a EMV para $\sigma^2$ es
        \begin{align*}
            \hat{\sigma}^2 = \frac{1}{n} \sum_{i=1}^{n} (X_i - \mu_0)^2.
        \end{align*}

        Sería necesario comprobar que este estadístico efectivamente es un estimador de $\sigma^2$. Como $\sigma^2 > 0$ y $\hat{\sigma}^2 > 0$, podemos concluir que efectivamente $\hat{\sigma}^2$ es un estimador de $\sigma^2$, y por lo tanto, $\hat{\sigma}^2$ es el EMV de $\sigma^2$.

        \item \ul{$\mu$ y $\sigma^2$ desconocidos:} Repidiendo el procedimiento anterior, se obtiene que los candidatos a EMV son
        \begin{align*}
            \hat{\mu} &= \frac{1}{n} \sum_{i=1}^{n} X_i = \ol{X}\\
            \hat{\sigma}^2 &= \frac{1}{n} \sum_{i=1}^{n} (X_i - \hat{\mu})^2 = \frac{1}{n} \sum_{i=1}^{n} (X_i - \ol{X})^2
        \end{align*}
    \end{enumerate}
\end{ejemplo}

\begin{ejemplo}
    Sea $X_1, \ldots, X_n$ una m.a.s. de $X\rightsquigarrow \cc{U}\{1, \ldots, N\}$, $N \in \mathbb{N}$. Encontrar el EMV para $N$.\\

    La función de verosimilitud es:
    \begin{align*}
        L_{x_1, \ldots, x_n}(N) &= \prod_{i=1}^{n} f_{X_i}(x_i; N) = \prod_{i=1}^{n} \frac{1}{N} = \frac{1}{N^n}\qquad \forall N \in \mathbb{N}, N \geq \max\{x_1, \ldots, x_n\} = x_{(n)}.
    \end{align*}

    Por tanto, podemos escribir la función de verosimilitud como:
    \begin{align*}
        L_{x_1, \ldots, x_n}(N) = \begin{cases}
            \dfrac{1}{N^n} & N \in \mathbb{N}, N \geq x_{(n)}\\ \\
            0 & N \in \mathbb{N}, N < x_{(n)}
        \end{cases}
    \end{align*}

    La función de verosimilitud no es continua, por lo que tampoco es diferenciable. Por tanto, no podemos derivar e igualar a cero para encontrar el máximo. Para encontrar el máximo, nos restringimos a $N\geq x_{(n)}$; $g=(L_{x_1, \ldots, x_n})_{|_{N\geq x_{(n)}}}$. Esta sí es derivable, luego:
    \begin{align*}
        g(N) = \frac{1}{N^n} \implies g'(N) = -\frac{n}{N^{n+1}} < 0 \qquad \forall N > x_{(n)}.
    \end{align*}

    Como además $g(x_{(n)})>0$, el máximo de la función de verosimilitud se alcanza en $N = x_{(n)}$, el candidato a EMV es:
    \begin{align*}
        \hat{N} = X_{(n)} = \max\{X_1, \ldots, X_n\}.
    \end{align*}

    Como $N\in \mathbb{N}$ y $\hat{N}\in \mathbb{N}$, podemos concluir que efectivamente $\hat{N}$ es un estimador de $N$, y por lo tanto, $\hat{N}$ es el EMV de $N$.
\end{ejemplo}

\begin{ejemplo}
    Sea $X_1, \ldots, X_n$ una m.a.s. de $X \rightsquigarrow \cc{U}(\theta - \nicefrac{1}{2}, \theta + \nicefrac{1}{2})$, $\theta \in \mathbb{R}$. Encontrar un EMV para $\theta$ y deducir que el EMV no tiene por qué ser único.\\

    La función de verosimilitud es:
    \begin{align*}
        L_{x_1, \ldots, x_n}(\theta) &= \prod_{i=1}^{n} f_{X_i}(x_i; \theta) = \prod_{i=1}^{n} 1 = 1 \qquad \forall \theta \in \mathbb{R}, \theta - \frac{1}{2} \leq x_i \leq \theta + \frac{1}{2}, i = 1, \ldots, n.
    \end{align*}

    Por tanto, considerando $x_{(1)} = \min\{x_1, \ldots, x_n\}$ y $x_{(n)} = \max\{x_1, \ldots, x_n\}$, podemos escribir la función de verosimilitud como:
    \begin{align*}
        L_{x_1, \ldots, x_n}(\theta) = \begin{cases}
            1 & \theta \in [x_{(n)} - \nicefrac{1}{2}, x_{(1)} + \nicefrac{1}{2}]\\
            0 & \theta \in \mathbb{R} \setminus [x_{(n)} - \nicefrac{1}{2}, x_{(1)} + \nicefrac{1}{2}]
        \end{cases}
    \end{align*}

    En este caso, no es necesario derivar para encontrar el máximo. Por tanto, el máximo de la función de verosimilitud se alcanza en cualquier $\theta \in [x_{(n)} - \nicefrac{1}{2}, x_{(1)} + \nicefrac{1}{2}]$, por lo que los candidatos a EMV son:
    \begin{align*}
        \hat{\theta} \in [X_{(n)} - \nicefrac{1}{2}, X_{(1)} + \nicefrac{1}{2}].
    \end{align*}

    Además, como $\theta\in \mathbb{R}$ y $\hat{\theta}\in \mathbb{R}$, podemos concluir que efectivamente $\hat{\theta}$ es un estimador de $\theta$, y por lo tanto, $\hat{\theta}$ es un EMV de $\theta$.
\end{ejemplo}

\begin{ejemplo}
    Sea $X$ una v.a. con distribución $\cc{B}(1, p)$ con $p \in \left[\nicefrac{1}{4}, \nicefrac{3}{4}\right]$. Calcular el EMV de $p$, para una muestra de tamaño 1, ver que no es insesgado y calcular su error cuadrático medio: $E[\hat{p} - p]^2$. Comprobar que el estimador $T(X) = \nicefrac{1}{2}$ es mejor que $\hat{p}$ en el sentido del ECM. (El EMV no tiene porque ser el mejor en el sentido del menor error cuadrático medio).\\

    La función de verosimilitud es:
    \begin{align*}
        L_{x_1, \ldots, x_n}(p) &= \prod_{i=1}^{n} f_{X_i}(x_i; p) = \prod_{i=1}^{n} p^{x_i}(1-p)^{1-x_i} = p^{\sum\limits_{i=1}^{n} x_i}(1-p)^{n - \sum\limits_{i=1}^{n} x_i}\qquad \forall p \in \left[\frac{1}{4}, \frac{3}{4}\right].
    \end{align*}

    Por tanto, la función de verosimilitud para una muestra de tamaño 1 es:
    \begin{align*}
        L_{x_1}(p) = p^{x_1}(1-p)^{1-x_1} \qquad \forall p \in \left[\frac{1}{4}, \frac{3}{4}\right].
    \end{align*}

    Definiéndola por partes, tenemos que:
    \begin{align*}
        L_{x_1}(p) = \begin{cases}
            p^{x_1}(1-p)^{1-x_1} & p \in \left[\nicefrac{1}{4}, \nicefrac{3}{4}\right]\\
            0 & p \in \mathbb{R} \setminus \left[\nicefrac{1}{4}, \nicefrac{3}{4}\right]
        \end{cases}
    \end{align*}

    Por tanto, para encontrar el máximo de la función de verosimilitud, nos restringimos a $p \in \left[\nicefrac{1}{4}, \nicefrac{3}{4}\right]$. Como $x_1 \in \{0, 1\}$, tenemos dos casos:
    \begin{enumerate}
        \item $x_1 = 0$: En este caso, la función de verosimilitud es $L_{0}(p) = 1 - p$, que es decreciente en $\left[\nicefrac{1}{4}, \nicefrac{3}{4}\right]$. Por tanto, el máximo se alcanza en $p = \nicefrac{1}{4}$, y el EMV es $\hat{p} = \nicefrac{1}{4}$.
        \item $x_1 = 1$: En este caso, la función de verosimilitud es $L_{1}(p) = p$, que es creciente en $\left[\nicefrac{1}{4}, \nicefrac{3}{4}\right]$. Por tanto, el máximo se alcanza en $p = \nicefrac{3}{4}$, y el EMV es $\hat{p} = \nicefrac{3}{4}$.
    \end{enumerate}

    Por tanto, el candidato a EMV es:
    \begin{align*}
        \hat{p} = \frac{1}{4}(1 - X) + \frac{3}{4}X = \frac{1}{4} + \frac{1}{2}X.
    \end{align*}

    Como $p \in \left[\nicefrac{1}{4}, \nicefrac{3}{4}\right]$ y $\hat{p} \in \left\{\nicefrac{1}{4}, \nicefrac{3}{4}\right\}$, podemos concluir que efectivamente $\hat{p}$ es un estimador de $p$, y por lo tanto, $\hat{p}$ es el EMV de $p$. Vemos que $\hat{p}$ no es insesgado, ya que:
    \begin{align*}
        E_p[\hat{p}] &= \frac{1}{4} + \frac{1}{2}E_p[X] = \frac{1}{4} + \frac{1}{2}p\qquad \forall p \in \left[\frac{1}{4}, \frac{3}{4}\right].
    \end{align*}

    Tenemos que:
    \begin{align*}
        E_p[\hat{p}] = p\iff \frac{1}{4} + \frac{1}{2}p = p \iff \frac{1}{4} = \frac{1}{2}p \iff p = \frac{1}{2}.
    \end{align*}

    Por tanto, como tendría que cumplirse para todo $p \in \left[\nicefrac{1}{4}, \nicefrac{3}{4}\right]$, podemos concluir que $\hat{p}$ no es insesgado. Calculamos el error cuadrático medio de $\hat{p}$:
    \begin{align*}
        E_p[(\hat{p} - p)^2] &= E_p\left[\left(\frac{1}{4} + \frac{1}{2}X - p\right)^2\right]
        = E_p\left[\left(\frac{1}{4} - p + \frac{1}{2}X\right)^2\right] \\
        &= \left(\frac{1}{4} - p\right)^2 + \frac{1}{4}E_p[X^2] + \left(\frac{1}{4} - p\right)E_p[X]
        =\\&= \left(\frac{1}{4} - p\right)^2 + \frac{1}{4}p + \left(\frac{1}{4} - p\right)p
        = \left(\frac{1}{4} - p\right)^2 + p\left(\frac{1}{2} - p\right)
        =\\&= \frac{1}{16} - \frac{1}{2}p + p^2 + \frac{1}{2}p - p^2 
        = \frac{1}{16} \qquad \forall p \in \left[\frac{1}{4}, \frac{3}{4}\right].
    \end{align*}

    Por otro lado, para el estimador $T(X) = \nicefrac{1}{2}$, tenemos:
    \begin{align*}
        E_p[(T(X) - p)^2] &= \left(\frac{1}{2} - p\right)^2.
    \end{align*}

    Dado que $p \in \left[\nicefrac{1}{4}, \nicefrac{3}{4}\right]$, se tiene que:
    \begin{align*}
        -\frac{1}{4} \le p - \frac{1}{2} \le \frac{1}{4} \implies \left|p - \frac{1}{2}\right| \le \frac{1}{4}.
    \end{align*}

    Elevando al cuadrado:
    \begin{align*}
        E_p[(T(X) - p)^2] = \left(p - \frac{1}{2}\right)^2 \le \left(\frac{1}{4}\right)^2 = \frac{1}{16} = E_p[(\hat{p} - p)^2] \qquad \forall p \in \left[\frac{1}{4}, \frac{3}{4}\right].
    \end{align*}
    alcanzándose la igualdad únicamente en los extremos $p = \nicefrac{1}{4}$ y $p = \nicefrac{3}{4}$, y siendo estrictamente positiva en todo el interior $\left]\nicefrac{1}{4}, \nicefrac{3}{4}\right[$. Por tanto, concluimos que $T(X)=~\nicefrac{1}{2}$ es un estimador estrictamente mejor que el EMV $\hat{p}$ en el interior del espacio paramétrico según el criterio del error cuadrático medio.
\end{ejemplo}

\begin{teo}[Relación entre EMV y estadístico suficiente]
    Sea $X$ una v.a. con distribución en una familia paramétrica de distribuciones, $\{F_\theta, \theta \in \Theta\}$. Sea $X_1, \ldots, X_n$ una m.a.s. de $X$. Supongamos que la familia admite un estadístico suficiente $T(X_1, \ldots, X_n)$. Entonces, si existe un EMV de $\theta$, es una función (no constante) de $T(X_1, \ldots, X_n)$.
\end{teo}
\begin{proof}
    Sea $T(X_1, \ldots, X_n)$ un estadístico suficiente para $\theta$. Por el Teorema de Factorización de Neyman-Fisher, la función de verosimilitud puede escribirse como:
    \begin{align*}
        L_{x_1, \ldots, x_n}(\theta) = f_{X_1, \ldots, X_n}(x_1, \ldots, x_n) = g_\theta(T(x_1, \ldots, x_n))h(x_1, \ldots, x_n) \qquad \forall \theta \in \Theta.
    \end{align*}
    Si existe un EMV $\hat{\theta}(X_1, \ldots, X_n)$ de $\theta$, entonces para cualquier $(x_1, \ldots, x_n) \in \mathbb{R}^n$:
    \begin{align*}
        \max_{\theta \in \Theta} L_{x_1, \ldots, x_n}(\theta) = L_{x_1, \ldots, x_n}(\hat{\theta}(x_1, \ldots, x_n)) = g_{\hat{\theta}(x_1, \ldots, x_n)}(T(x_1, \ldots, x_n))h(x_1, \ldots, x_n)
    \end{align*}
    
    Como $h(x_1, \ldots, x_n)$ no depende de $\theta$, maximizar $L_{x_1, \ldots, x_n}(\theta)$ respecto a $\theta$ es equivalente a maximizar $g_\theta(T(x_1, \ldots, x_n))$ respecto a $\theta$. Por tanto, para cualquier $(x_1, \ldots, x_n) \in \mathbb{R}^n$, se tiene que:
    \begin{align*}
        \max_{\theta \in \Theta} g_\theta(T(x_1, \ldots, x_n)) = g_{\hat{\theta}(x_1, \ldots, x_n)}(T(x_1, \ldots, x_n)).
    \end{align*}

    Por tanto, $\hat{\theta}(x_1, \ldots, x_n)$ es el valor de $\theta$ que maximiza $g_\theta(T(x_1, \ldots, x_n))$, y por tanto, $\hat{\theta}(x_1, \ldots, x_n)$ es una función (no constante) de $T(x_1, \ldots, x_n)$.    
\end{proof}

\begin{ejemplo}
    Sea $X_1, \ldots, X_n$ una m.a.s. de $X \rightsquigarrow \cc{U}(\theta, \theta + 1)$, $\theta \in \mathbb{R}$. Calcular el estadístico suficiente, un EMV para $\theta$, y comprobar que no coinciden.\\

    En vistas de aplicar el Teorema de Factorización de Neyman-Fisher, calculamos la función de densidad conjunta:
    \begin{align*}
        f_{\theta}^n(x_1, \ldots, x_n) &= \prod_{i=1}^{n} f_{\theta}(x_i) = \prod_{i=1}^{n} 1 = 1 \qquad \forall x_i \in [\theta, \theta + 1], i = 1, \ldots, n.
    \end{align*}

    Por tanto, considerando $x_{(1)} = \min\{x_1, \ldots, x_n\}$ y $x_{(n)} = \max\{x_1, \ldots, x_n\}$, podemos escribir la función de densidad conjunta como:
    \begin{align*}
        f_{\theta}^n(x_1, \ldots, x_n) = \ind_{\bb{R}^+}(x_{(1)}-\theta)\ind_{\bb{R}^+}(\theta + 1 - x_{(n)}) \qquad \forall \theta \in \mathbb{R}.
    \end{align*}

    Por el Teorema de Factorización de Neyman-Fisher, el estadístico suficiente es:
    \begin{align*}
        T(X_1, \ldots, X_n) = (X_{(1)}, X_{(n)}) = (\min\{X_1, \ldots, X_n\}, \max\{X_1, \ldots, X_n\}).
    \end{align*}

    Por otro lado, la función de verosimilitud es:
    \begin{align*}
        L_{x_1, \ldots, x_n}(\theta) &= \prod_{i=1}^{n} f_{\theta}(x_i) = \prod_{i=1}^{n} 1 = 1 \qquad \forall \theta \in [x_{(n)} - 1, x_{(1)}].
    \end{align*}

    Por tanto, el máximo de la función de verosimilitud se alcanza en cualquier $\theta \in [x_{(n)} - 1, x_{(1)}]$, por lo que los candidatos a EMV son:
    \begin{align*}
        \hat{\theta} \in [X_{(n)} - 1, X_{(1)}].
    \end{align*}

    En particular, podemos tomar $\hat{\theta} = X_{(1)}$ o $\hat{\theta} = X_{(n)} - 1$. Concluimos que el EMV no coincide con el estadístico suficiente, ya que $T(X_1, \ldots, X_n) = (X_{(1)}, X_{(n)})$ y $\hat{\theta} \in [X_{(n)} - 1, X_{(1)}]$. No obstante, sí es una función (no constante) del estadístico suficiente.
\end{ejemplo}

\begin{teo}[Relación entre EMV y estimador eficiente]
    Sea $X$ una v.a. con distribución en una familia paramétrica de distribuciones, $\{F_\theta, \theta \in \Theta\}$, regular con $0 < I_X(\theta) < \infty$. Sea $X_1, \ldots, X_n$ una m.a.s. de $X$. Si $T(X_1, \ldots, X_n)$ es un estimador eficiente para $\theta$, entonces existe un único EMV y coincide con $T(X_1, \ldots, X_n)$.
    \begin{proof}
        Sea $T(X_1, \ldots, X_n)$ un estimador eficiente para $\theta$. Por la caracterización de los estimadores eficientes, para todo $\theta \in \Theta$ existe $a(\theta)\neq 0$ tal que:
        \begin{align*}
            \frac{\partial}{\partial \theta} \ln(L_{X_1, \ldots, X_n}(\theta)) = a(\theta)(T(X_1, \ldots, X_n) - \theta).
        \end{align*}

        Consideramos la función de verosimilitud $L_{x_1, \ldots, x_n}(\theta)$ para una muestra $(x_1, \ldots, x_n) \in \mathbb{R}^n$:
        \begin{align*}
            L_{x_1, \ldots, x_n}(\theta) = f_{X_1, \ldots, X_n}(x_1, \ldots, x_n) \qquad \forall \theta \in \Theta.
        \end{align*}

        Como el logaritmo es una función estrictamente creciente, podemos trabajar con la log-verosimilitud:
        \begin{align*}
            \ln(L_{x_1, \ldots, x_n}(\theta)) = \ln(f_{X_1, \ldots, X_n}(x_1, \ldots, x_n)) \qquad \forall \theta \in \Theta.
        \end{align*}

        Como se trata de buscar el máximo de la función de verosimilitud, derivamos e igualamos a cero:
        \begin{align*}
            \frac{\partial}{\partial \theta} \ln(L_{x_1, \ldots, x_n}(\theta)) = \frac{\partial}{\partial \theta} \ln(f_{X_1, \ldots, X_n}(x_1, \ldots, x_n)) = a(\theta)(T(x_1, \ldots, x_n) - \theta) = 0.
        \end{align*}

        Como $a(\theta)\neq 0$, se tiene que:
        \begin{align*}
            T(x_1, \ldots, x_n) - \theta = 0 \implies \theta = T(x_1, \ldots, x_n).
        \end{align*}

        Por tanto, el único candidato a EMV es:
        \begin{align*}
            \hat{\theta} = T(X_1, \ldots, X_n).
        \end{align*}

        Como $T(X_1, \ldots, X_n)$ es un estimador de $\theta$, podemos concluir que efectivamente $\hat{\theta}$ es un estimador de $\theta$, y por lo tanto, $\hat{\theta}$ es el único EMV de $\theta$.
    \end{proof}
\end{teo}

\begin{teo}[Invariancia de Zehna]
    Sea X una v.a. con distribución en una familia paramétrica de distribuciones, $\{F_\theta, \theta \in \Theta\}$. Sea $X_1, \ldots, X_n$ una m.a.s. de $X$. Sea $g$ una función medible. Si $\hat{\theta}(X_1, \ldots, X_n)$ es EMV de $\theta$, entonces $g(\hat{\theta}(X_1, \ldots, X_n))$ es EMV de $g(\theta)$.
    \begin{proof}
        Sea $\hat{\theta}(X_1, \ldots, X_n)$ un EMV de $\theta$. Por definición, para cualquier $(x_1, \ldots, x_n) \in \mathbb{R}^n$:
        \begin{align*}
            \max_{\theta \in \Theta} L_{x_1, \ldots, x_n}(\theta) = L_{x_1, \ldots, x_n}(\hat{\theta}(x_1, \ldots, x_n)).
        \end{align*}

        Veamos que $g(\hat{\theta}(X_1, \ldots, X_n))$ es EMV de $g(\theta)$. Para ello, consideramos la función de verosimilitud respecto a $g(\theta)$:
        \begin{align*}
            M_{x_1, \ldots, x_n}(g(\theta)) = \sup_{\theta^* \in \Theta : g(\theta^*) = g(\theta)} L_{x_1, \ldots, x_n}(\theta^*) \qquad \forall g(\theta) \in g(\Theta).
        \end{align*}

        Queremos ver que:
        \begin{equation*}
            M_{x_1, \ldots, x_n}(g(\hat{\theta}(x_1, \ldots, x_n))) = \max_{g(\theta) \in g(\Theta)} M_{x_1, \ldots, x_n}(g(\theta)).
        \end{equation*}

        Tenemos que:
        \begin{align*}
            M_{x_1, \ldots, x_n}(g(\hat{\theta}(x_1, \ldots, x_n))) &= \sup_{\theta^* \in \Theta : g(\theta^*) = g(\hat{\theta}(x_1, \ldots, x_n))} L_{x_1, \ldots, x_n}(\theta^*) \\
            &= L_{x_1, \ldots, x_n}(\hat{\theta}(x_1, \ldots, x_n)) \\
            &= \max_{\theta \in \Theta} L_{x_1, \ldots, x_n}(\theta) \\
            &\geq \sup_{\theta^* \in \Theta : g(\theta^*) = g(\theta)} L_{x_1, \ldots, x_n}(\theta^*) = M_{x_1, \ldots, x_n}(g(\theta)) \qquad \forall g(\theta) \in g(\Theta).
        \end{align*}

        Por tanto, $g(\hat{\theta}(X_1, \ldots, X_n))$ es EMV de $g(\theta)$.
    \end{proof}
\end{teo}


\begin{ejemplo}
    La duración de cierto tipo de lámparas sigue una distribución exponencial con media $\theta$, desconocida. Después de observar el tiempo de vida de $n$ lámparas, estimar por máxima verosimilitud la probabilidad de que la duración de una lámpara sea superior a $500$ horas.\\

    Sea $X$ una variable aleatoria que representa la duración de una lámpara, con distribución $X \rightsquigarrow \exp(a)$, $a > 0$. Sabiendo la que la esperanza de $X$ es $\theta$, tenemos que:
    \begin{align*}
        E[X] = \theta = \frac{1}{a} \implies a = \frac{1}{\theta}.
    \end{align*}

    Por tanto, $X \rightsquigarrow \exp(\nicefrac{1}{\theta})$, $\theta > 0$. Buscamos el EMV de $P_{\theta}[X > 500]$. En primer lugar, calculamos:
    \begin{align*}
        P_{\theta}[X > 500] = \int_{500}^{\infty} \frac{1}{\theta} \exp\left(-\frac{x}{\theta}\right) dx = \left[-\exp\left(-\frac{x}{\theta}\right)\right]_{500}^{\infty} = \exp\left(-\frac{500}{\theta}\right)
    \end{align*}

    Por tanto, buscamos el EMV de $g(\theta) = P_{\theta}[X > 500] = \exp(-\nicefrac{500}{\theta})$. Por el Teorema de Invariancia de Zehna, si $\hat{\theta}$ es el EMV de $\theta$, entonces:
    \begin{align*}
        g(\hat{\theta}) = \exp\left(-\frac{500}{\hat{\theta}}\right)
    \end{align*}
    es el EMV de $g(\theta) = P_{\theta}[X > 500]$. Por tanto, nos centramos en calcular el EMV de $\theta$. La función de verosimilitud es:
    \begin{align*}
        L_{x_1, \ldots, x_n}(\theta) &= \prod_{i=1}^{n} f_{X_i}(x_i; \theta) = \prod_{i=1}^{n} \frac{1}{\theta} \exp\left(-\frac{x_i}{\theta}\right) = \frac{1}{\theta^n} \exp\left(-\frac{1}{\theta} \sum_{i=1}^{n} x_i\right) \qquad \forall \theta > 0
    \end{align*}

    Puesto que el logaritmo es una función estrictamente creciente, podemos trabajar con la log-verosimilitud:
    \begin{align*}
        \ln(L_{x_1, \ldots, x_n}(\theta)) &= -n\ln(\theta) - \frac{1}{\theta} \sum_{i=1}^{n} x_i \qquad \forall \theta > 0
    \end{align*}

    Como buscamos el máximo, derivamos e igualamos a cero:
    \begin{align*}
        \frac{d}{d\theta} \ln(L_{x_1, \ldots, x_n}(\theta)) &= -\frac{n}{\theta} + \frac{1}{\theta^2} \sum_{i=1}^{n} x_i = 0 \implies -n\theta + \sum_{i=1}^{n} x_i = 0 \implies \theta = \frac{1}{n} \sum_{i=1}^{n} x_i
    \end{align*}

    Sería necesario comprobar que efectivamente se trata de un máximo, pero se tiene. Por tanto, el candidato a EMV de $\theta$ es:
    \begin{align*}
        \hat{\theta} = \frac{1}{n} \sum_{i=1}^{n} X_i = \ol{X}.
    \end{align*}

    Como $\theta > 0$ y $\hat{\theta} > 0$, podemos concluir que efectivamente $\hat{\theta}$ es un estimador de $\theta$, y por lo tanto, $\hat{\theta}$ es el EMV de $\theta$. Por el Teorema de Invariancia de Zehna, si $\hat{\theta}$ es el EMV de $\theta$, entonces:
    \begin{align*}
        g(\hat{\theta}) = \exp\left(-\frac{500}{\hat{\theta}}\right)
    \end{align*}
    es el EMV de $g(\theta) = P_{\theta}[X > 500]$. Por tanto, el EMV de $P_{\theta}[X > 500]$ es:
    \begin{align*}
        \widehat{P_{\theta}[X > 500]} = \exp\left(-\frac{500}{\hat{\theta}}\right) = \exp\left(-\frac{500}{\ol{X}}\right).
    \end{align*}
\end{ejemplo}

\begin{ejemplo}
    En el muestreo de una v.a. $X \rightsquigarrow \cc{P}(\lambda)$, $\lambda > 0$, se obtiene que en $n$ observaciones aparece $y$ veces el valor $0$. Obtener un EMV de $\lambda$ a partir de esta información.\\

    Como primer enfoque, buscaríamos definir la función de verosimilitud de $\lambda$, $L_{x_1, \ldots, x_n}(\lambda)$. No obstante, como no conocemos los valores de las observaciones, no podemos seguir este enfoque. Definimos la variable auxiliar $Y$ como el número de veces que aparece el valor $0$ en la muestra, luego:
    \begin{align*}
        Y = \sum_{i=1}^{n} I_{\{0\}}(X_i) \rightsquigarrow \cc{B}(n, P[X = 0]) = \cc{B}(n, e^{-\lambda}).
    \end{align*}

    Por tanto, la función de verosimilitud de $e^{-\lambda}$ es:
    \begin{align*}
        L_{y}(e^{-\lambda}) &= P[Y = y] = \binom{n}{y} (e^{-\lambda})^y (1 - e^{-\lambda})^{n - y} \qquad \forall e^{-\lambda} \in \left]0, 1\right[
    \end{align*}

    Puesto que el logaritmo es una función estrictamente creciente, podemos trabajar con la log-verosimilitud:
    \begin{align*}
        \ln(L_{y}(e^{-\lambda})) &= \ln\left(\binom{n}{y}\right) + y\ln(e^{-\lambda}) + (n - y)\ln(1 - e^{-\lambda}) \qquad \forall e^{-\lambda} \in \left]0, 1\right[
    \end{align*}

    Como buscamos el máximo, derivamos e igualamos a cero:
    \begin{align*}
        \frac{d}{d(e^{-\lambda})} \ln(L_{y}(e^{-\lambda})) &= \frac{y}{e^{-\lambda}} - \frac{n - y}{1 - e^{-\lambda}} = 0 \implies \\&\implies y(1 - e^{-\lambda}) - (n - y)e^{-\lambda} = 0 \implies \\&\implies y - ne^{-\lambda} = 0 \implies e^{-\lambda} = \frac{y}{n}
    \end{align*}

    Sería necesario comprobar que efectivamente se trata de un máximo, pero se tiene. Por tanto, el candidato a EMV de $e^{-\lambda}$ es:
    \begin{align*}
        \widehat{e^{-\lambda}} = \frac{Y}{n}.
    \end{align*}

    Como $e^{-\lambda} \in \left]0, 1\right[$ y $\widehat{e^{-\lambda}} \in \left[0, 1\right]$, podemos concluir que efectivamente $\widehat{e^{-\lambda}}$ es un estimador de $e^{-\lambda}$, y por lo tanto, $\widehat{e^{-\lambda}}$ es el EMV de $e^{-\lambda}$. Por el Teorema de Invariancia de Zehna, si $\widehat{\lambda}$ es el EMV de $\lambda$, entonces:
    \begin{align*}
        e^{-\widehat{\lambda}} = \widehat{e^{-\lambda}} \implies \widehat{\lambda} = -\ln(\widehat{e^{-\lambda}}) = -\ln\left(\frac{Y}{n}\right) = \ln\left(\frac{n}{Y}\right).
    \end{align*}
\end{ejemplo}

\begin{ejemplo}
    Estimar mediante el método de los momentos los parámetros de las siguientes distribuciones:
    \begin{enumerate}
        \item $X \rightsquigarrow \cc{N}(\mu, \sigma^2)$.
        
        Puesto que tenemos dos parámetros, necesitamos dos momentos. Por tanto, calculamos el primer y segundo momento poblacionales:
        \begin{align*}
            m_1 &= E[X] = \mu\\
            m_2 &= E[X^2] = Var[X] + (E[X])^2 = \sigma^2 + \mu^2
        \end{align*}

        Por otro lado, calculamos el primer y segundo momento muestrales:
        \begin{align*}
            A_1 &= \frac{1}{n} \sum_{i=1}^{n} X_i = \ol{X}\\
            A_2 &= \frac{1}{n} \sum_{i=1}^{n} X_i^2
        \end{align*}

        Por el método de los momentos, igualamos los momentos poblacionales con los muestrales:
        \begin{align*}
            m_1 &= A_1 \implies \mu = \ol{X}\\
            m_2 &= A_2 \implies \sigma^2 + \mu^2 = A_2 \implies \sigma^2 = \frac{1}{n} \sum_{i=1}^{n} X_i^2 - \ol{X}^2 = Var[X_1, \ldots, X_n]
        \end{align*}

        Por tanto, los estimadores por el método de los momentos son:
        \begin{align*}
            \hat{\mu} = \ol{X},\qquad \hat{\sigma}^2 = Var[X_1, \ldots, X_n]
        \end{align*}

        Sería necesario comprobar que efectivamente $\hat{\mu}$ y $\hat{\sigma}^2$ son estimadores de $\mu$ y $\sigma^2$, respectivamente, pero se tiene. Por tanto, podemos concluir que $\hat{\mu}$ y $\hat{\sigma}^2$ son los estimadores por el método de los momentos de $\mu$ y $\sigma^2$, respectivamente. Como podemos ver, los estimadores por el método de los momentos no tienen por qué ser insesgados.

        \item $X \rightsquigarrow \cc{U}(a, b)$.
        
        Puesto que tenemos dos parámetros, necesitamos dos momentos. Por tanto, calculamos el primer y segundo momento poblacionales:
        \begin{align*}
            m_1 &= E[X] = \frac{a + b}{2}\\
            m_2 &= E[X^2] = Var[X] + (E[X])^2 = \frac{(b - a)^2}{12} + \left(\frac{a + b}{2}\right)^2 =\\&= \frac{b^2 +a^2-2ab + 3a^2 + 3b^2 + 6ab}{12} = \frac{a^2 + ab + b^2}{3}
        \end{align*}

        Por otro lado, calculamos el primer y segundo momento muestrales:
        \begin{align*}
            A_1 &= \frac{1}{n} \sum_{i=1}^{n} X_i = \ol{X}\\
            A_2 &= \frac{1}{n} \sum_{i=1}^{n} X_i^2
        \end{align*}

        Por el método de los momentos, igualamos los momentos poblacionales con los muestrales:
        \begin{align*}
            m_1 &= A_1 \implies \frac{a + b}{2} = \ol{X} \implies a + b = 2\ol{X} \\
            m_2 &= A_2 \implies \frac{a^2 + ab + b^2}{3} = A_2
        \end{align*}

        Resolviendo el sistema de ecuaciones, obtenemos:
        \begin{align*}
            a &= \ol{X} - \sqrt{3}\sqrt{A_2 - \ol{X}^2}
            b &= \ol{X} + \sqrt{3}\sqrt{A_2 - \ol{X}^2}
        \end{align*}

        Por tanto, los estimadores por el método de los momentos son:
        \begin{align*}
            \hat{a} = \ol{X} - \sqrt{3}\sqrt{A_2 - \ol{X}^2},\qquad \hat{b} = \ol{X} + \sqrt{3}\sqrt{A_2 - \ol{X}^2}
        \end{align*}

        Sería necesario comprobar que efectivamente $\hat{a}$ y $\hat{b}$ son estimadores de $a$ y $b$, respectivamente, pero se tiene. Por tanto, podemos concluir que $\hat{a}$ y $\hat{b}$ son los estimadores por el método de los momentos de $a$ y $b$, respectivamente.

        \item $X \rightsquigarrow \cc{P}(\lambda)$.
        
        Puesto que tenemos un parámetro, necesitamos un momento. Por tanto, calculamos el primer momento poblacional:
        \begin{align*}
            m_1 &= E[X] = \lambda
        \end{align*}

        Por otro lado, calculamos el primer momento muestral:
        \begin{align*}
            A_1 &= \frac{1}{n} \sum_{i=1}^{n} X_i = \ol{X}
        \end{align*}

        Por el método de los momentos, igualamos los momentos poblacionales con los muestrales:
        \begin{align*}
            m_1 &= A_1 \implies \lambda = \ol{X}
        \end{align*}

        Por tanto, el estimador por el método de los momentos es:
        \begin{align*}
            \hat{\lambda} = \ol{X}
        \end{align*}

        \begin{observacion}
            Notemos que, si hubiésemos incluido las igualdades relacionadas con el segundo momento, habríamos obtenido $\hat{\lambda} = Var[X_1, \ldots, X_n]$, por lo que es importante añadir tan solo las igualdades necesarias para obtener el número de ecuaciones igual al número de parámetros.
        \end{observacion}
        \item $X \rightsquigarrow \cc{U}(0, \theta)$.
        
        Puesto que tenemos un parámetro, necesitamos un momento. Por tanto, calculamos el primer momento poblacional:
        \begin{align*}
            m_1 &= E[X] = \frac{\theta}{2}
        \end{align*}

        Por otro lado, calculamos el primer momento muestral:
        \begin{align*}
            A_1 &= \frac{1}{n} \sum_{i=1}^{n} X_i = \ol{X}
        \end{align*}

        Por el método de los momentos, igualamos los momentos poblacionales con los muestrales:
        \begin{align*}
            m_1 &= A_1 \implies \frac{\theta}{2} = \ol{X} \implies \theta = 2\ol{X}
        \end{align*}

        Por tanto, el estimador por el método de los momentos es:
        \begin{align*}
            \hat{\theta} = 2\ol{X}
        \end{align*}

        Sería necesario comprobar que efectivamente $\hat{\theta}$ es un estimador de $\theta$, pero se tiene. Por tanto, podemos concluir que $\hat{\theta}$ es el estimador por el método de los momentos de $\theta$.

        \begin{observacion}
            Veamos qué ocurre si la muestra que consideramos es $(2,4,9,1)$. En este caso:
            \begin{align*}
                \hat{\theta} = 2\ol{X} = 2\cdot \frac{2 + 4 + 9 + 1}{4} = 2\cdot \frac{16}{4} = 8
            \end{align*}

            Por tanto, el estimador por el método de los momentos nos da $\hat{\theta} = 8$, luego podríamos pensar que $X\rightsquigarrow \cc{U}(0, 8)$. No obstante, esto no es posible, ya que la muestra contiene el valor $9$, que no pertenece al intervalo $[0, 8]$. Por tanto, el estimador por el método de los momentos no siempre es adecuado.
        \end{observacion}
    \end{enumerate}
\end{ejemplo}

\begin{ejemplo}
    Para estmar la aceleración de la gravedad en una ciudad, $\theta$, se deja caer un objeto durante tiempos $t_1, \ldots, t_n$ y se mide el espacio recorrido en cada tiempo. Si $X_i$ representa la medida correspondiente al espacio recorrido en el tiempo $t_i$, con error de medida $\varepsilon_i$, estimar $\theta$ por el método de mínimos cuadrados.
    
    Por la ecuación del movimiento rectilíneo uniformemente acelerado, tenemos que:
    \begin{align*}
        X_i = \frac{1}{2}\theta t_i^2 + \varepsilon_i \implies \varepsilon_i = X_i - \frac{1}{2}\theta t_i^2
    \end{align*}

    Por el método de mínimos cuadrados, buscamos minimizar la suma de los cuadrados de los errores:
    \begin{align*}
        S(\theta) = \sum_{i=1}^{n} \varepsilon_i^2 = \sum_{i=1}^{n} \left(x_i - \frac{1}{2}\theta t_i^2\right)^2
    \end{align*}

    Como buscamos el mínimo, derivamos e igualamos a cero:
    \begin{align*}
        \frac{d}{d\theta} S(\theta) &= \sum_{i=1}^{n} 2\left(x_i - \frac{1}{2}\theta t_i^2\right)\left(-\frac{1}{2}t_i^2\right) = 0 \implies \sum_{i=1}^{n} t_i^2 x_i - \frac{1}{2}\theta \sum_{i=1}^{n} t_i^4 = 0
    \end{align*}

    Por tanto, obtenemos:
    \begin{align*}
        \theta = \frac{2\sum\limits_{i=1}^{n} t_i^2 x_i}{\sum\limits_{i=1}^{n} t_i^4}
    \end{align*}

    Por tanto, el estimador por el método de mínimos cuadrados es:
    \begin{align*}
        \hat{\theta} = \frac{2\sum\limits_{i=1}^{n} t_i^2 X_i}{\sum\limits_{i=1}^{n} t_i^4}
    \end{align*}


\end{ejemplo}